
\section{Related Literature}
Illicit financial activity detection has been an active area of research for decades, encompassing a wide range of applications such as financial fraud detection~\parencite{fan2025unearthing, yiu2014deterrence, wu2024fair}, anti–money laundering detection~\parencite{deprez2025network}, e-commerce fraud detection~\parencite{nanduri2020microsoft, luca2016fake}, social security fraud detection~\parencite{van2017gotcha}, securities fraud detection in emerging markets~\parencite{guo2022hierarchical}, etc. With the growing scale and complexity of financial transactions, machine learning has become a dominant approach for detecting illicit activities, offering improved accuracy and adaptability compared to traditional rule-based systems~\parencite{august2025impact, fan2025unearthing, cecchini2010detecting, xiao2023black}. This study focuses specifically on the challenges of anti–money laundering detection.

% related work on detection model
The traditional ML-based AML systems are poorly equipped to handle this environment~\parencite{pettersson2022combating, alotibi2022money, jensen2023fighting, lokanan2024predicting}. Eric and Jannis \parencite*{pettersson2022combating} evaluate four classic supervised learning algorithms (e.g., logistic regression, random forest, support vector machines, and decision trees) on the AML task and find that these models are too slow and require further optimization for the anti–money laundering context. More advanced supervised machine learning models have been proposed, trained on historical transaction data to detect anomalies and laundering patterns, offering more accurate and adaptive detection than traditional rule-based systems~\parencite{golightly2023securing}. However, these methods rely on static machine learning and deep learning models, which fail to integrate information across institutions and struggle to adapt to temporally complex laundering strategies. Furthermore, studies show that some laundering schemes are executed rapidly over a handful of transactions, while others evolve over months and mix tactics such as layering and cross-border transfers~\parencite{amoako2025exploring}. Anti–money laundering detection models tuned to one temporal horizon often miss others, creating a persistent tension between premature intervention and delayed action. Wang et al. \parencite*{wang2024temporal} demonstrate that traditional anti–money laundering systems struggle to detect complex schemes that leverage cross-jurisdictional gaps and the temporal dispersion of transactions, thus failing to capture complex evolving patterns in transaction networks. 

% related work on graph-based methods for financial fraud and AML
Recently, Kurshan and Shen~\parencite*{kurshan2020graph} and Lokanan~\parencite*{lokanan2023financial} show the prominent performance of graph-based approaches on general financial fraud. For instance, Kurshan and Shen~\parencite*{kurshan2020graph} provide a high-level discussion of the challenges faced by machine learning methods in combating financial crime, emphasizing issues such as scalability, data quality, and operational constraints. Although their survey does not center on AML, they underscore the strong potential of network-based methods for modeling the large-scale transactional data handled by financial institutions. Complementary to these methodological surveys, Gerbrands et al.~\parencite*{gerbrands2022effect} examine AML from a policy and network-structure perspective, demonstrating how regulatory interventions reshape transactional networks and influence illicit financial behaviors. While this line of work highlights the importance of network analysis for understanding money-laundering dynamics, it does not develop learning-based detection models. More closely aligned with graph neural network approaches for AML, Kute et al.~\parencite*{kute2022explainable, kute2024explainable} demonstrate how deep graph learning models can be augmented with explanation mechanisms to clarify which transaction patterns and behavioral attributes drive model outputs. 


% related work on interpretable methods for AML
Model interpretability is critical in financial applications, as compliance officers and investigators must understand, trust, and act upon the model’s recommendations. Contemporary work on interpretable AML detection spans methodological advances in feature-based models, post-hoc explanation techniques, and network-aware methods and much of this literature appears in business-oriented outlets. Early applied studies established the value of supervised learning for prioritizing alerts while emphasizing model transparency~\parencite{chen2018machine, han2020artificial}. Building on these works, Jullum et al. \parencite*{jullum2020detecting} develop and validate machine-learning scoring systems for transaction prioritization, showing how transparent feature design and risk-scoring improve investigatory efficiency. Konstantinidis and Gegov \parencite*{konstantinidis2024deep} evaluate explainable AI tools (e.g., SHAP and LIME) as operational supplements to black-box models, demonstrating that feature-attribution explanations improve auditability and the quality of Suspicious Activity Reports. Oloyede \parencite*{oloyede2025evaluating} explores how interpretability (e.g., transparent ML models) affects regulatory compliance, trust, and operational effectiveness in financial institutions.  In parallel, Gu, Yang, and Liu \parencite*{gu2025optimization} incorporates causal reasoning to improve both the transparency and the decision consistency of AML systems. Recent reviews~\parencite{mazumder2025explainable} synthesize these strands and call explicitly for XAI and socio-technical integration so that automated detection maps cleanly onto operational workflows and regulatory expectations. Collectively, these works show that effective AML solutions must go beyond predictive accuracy to provide clear, auditable rationales for their decisions, thereby enhancing investigator trust, supporting compliance verification, and ensuring that automated detection remains aligned with financial governance standards.

% related work on decision/policy-making methods for AML
In the past decades, a growing body of research examines how anti–money laundering decisions and policies are shaped by regulatory frameworks, organizational practices, and behavioral responses to financial crime risks. Early work \parencite{ross2007money} highlights the central tension in AML supervision between rule-based requirements and risk-based decision-making, showing how compliance officers must interpret ambiguous regulations under conditions of limited information. Unger and Van Waarden \parencite*{unger2009dodge} further compare rule-based and risk-based AML strategies, arguing that excessive data collection without effective prioritization can overwhelm institutions and undermine detection quality. Complementing these policy-oriented perspectives, McCarthy, van Santen, and Fiedler \parencite*{mccarthy2015modeling} develop a microtheoretical model of money launderer behavior, illustrating how adversarial adaptation should inform AML policy design. Similarly, research on strategic learning behavior in markets demonstrates that sophisticated actors rapidly adapt their strategies in response to observed outcomes and platform policies~(\cite{xu2022ni}), underscoring the importance of adaptive detection systems that evolve alongside increasingly sophisticated laundering tactics. Helgesson and Mörth \parencite*{helgesson2016involuntary} shift attention to the role of professional intermediaries, demonstrating that lawyers become de facto policy actors when implementing AML and counter-terrorism financing mandates, often shaping regulatory outcomes through everyday compliance decisions. More recently, methodological advances have emerged, such as the reinforcement learning framework proposed by \parencite{rao2025reinforcement}, which models cross-border transaction anomalies from a behavioral-economics perspective and offers a data-driven approach to dynamic AML decision-making. Together, these studies reveal an evolving landscape in which AML effectiveness depends not only on formal rules and technologies but also on how institutions interpret risks, allocate attention, and adapt to strategic adversaries.

Despite significant progress, these AML detection methods suffer from several critical limitations that prevent them from keeping pace with increasingly complex laundering behaviors. Traditional supervised machine learning and deep learning models~\parencite{pettersson2022combating, alotibi2022money, jensen2023fighting, lokanan2024predicting, chen2018machine}, while effective in narrow settings, remain largely static, siloed, and institution-specific. These systems struggle to generalize across heterogeneous financial institutions, cannot integrate cross-bank intelligence, and often fail under data imbalance or rapidly evolving laundering tactics. Recent interpretable or explainable AI approaches~\parencite{chen2018machine, han2020artificial, kute2024explainable, kute2022explainable} focus primarily on feature-level transparency and do not reveal coordinated group behaviors or network structures underlying illicit flows. At the policy and decision-making level, current systems~\parencite{ross2007money, unger2009dodge, helgesson2016involuntary, mccarthy2015modeling} are still reactive, offering little guidance on when and how local institutions should intervene in unfolding laundering schemes under a global coordination system. As a result, AML operations remain fragmented: models detect anomalies, investigators manually infer group structures, and compliance teams determine actions without principled support. The privacy-preserving adaptive framework proposed in this paper directly addresses these gaps. By enabling graph-based federated learning, our method preserves institutional privacy while facilitating cross-bank sharing. The integration of focal-loss–based detection, cross-bank PPR based group identification, and hierarchical reinforcement learning creates a unified pipeline that not only identifies suspicious entities but also uncovers their coordinated structures and recommends cost-sensitive, context-aware interventions. In doing so, the proposed system advances AML practice from isolated, static prediction tools toward proactive, explainable, and strategically optimized financial crime prevention.


